% ═══════════════════════════════════════════════════════════════════════
\section{Remise et ressources}
% ═══════════════════════════════════════════════════════════════════════

\begin{frame}{Avant de remettre}
  \begin{Rem}{Liste de vérification finale}
    \small
    \begin{enumerate}
      \item Les deux documents compilent sans erreur (aucun message rouge).
      \item L'activité fait 2 à 3 pages et contient toutes les rubriques
        demandées, avec au moins une figure ou un tableau.
      \item L'examen affiche la grille et le bon total ; le corrigé
        s'obtient en activant \cmd{printanswers}.
      \item Tout est en français ; les décimaux ont une virgule.
      \item Les mathématiques ont été vérifiées par vous, pas seulement
        par l'IA.
      \item Vous pouvez expliquer à quoi sert chaque ligne du préambule.
    \end{enumerate}
    À remettre : \texttt{activite.pdf}, \texttt{examen.pdf},
    \texttt{examen-corrige.pdf}, et les fichiers \texttt{.tex}.
  \end{Rem}
\end{frame}

\begin{frame}{Pour aller plus loin}
  \small
  \begin{itemize}
    \item Overleaf, Learn LaTeX in 30 minutes et la documentation en ligne :
      \texttt{overleaf.com/learn}.
    \item Detexify, pour trouver un symbole en le dessinant :
      \texttt{detexify.kirelabs.org}.
    \item La documentation de la classe exam, très complète, avec de
      nombreux exemples : \texttt{ctan.org/pkg/exam}.
    \item La documentation de siunitx : \texttt{ctan.org/pkg/siunitx}.
    \item Les trois fichiers d'exemple de ce laboratoire :
      \texttt{01-minimal.tex}, \texttt{02-activite.tex},
      \texttt{03-examen.tex}.
  \end{itemize}
  \vspace{6pt}
  \begin{Classe}{Pour votre pratique}
    \small Commencez une banque de questions dès vos premiers stages : un fichier
    \texttt{exam} par notion, que vous enrichirez et réutiliserez pendant
    toute votre carrière.
  \end{Classe}
\end{frame}

\appendix

\begin{frame}[fragile]{Solutions de l'exercice « Reproduire ces formules »}
\begin{CodeTeX}[title={Une solution possible}]
% 1.
$f(x) = -2(x - 3)^2 + 1$
% 2.
L'aire d'un disque de rayon \qty{5}{\centi\metre} est
$A = \pi r^2 \approx \qty{78,54}{\centi\metre\squared}$.
% 3.
\[ x = \frac{-b \pm \sqrt{b^2 - 4ac}}{2a} \]
% 4.
\begin{align*}
  2(x + 4) &= 18 \\
  x + 4 &= 9 \\
  x &= 5
\end{align*}
% 5.
$\sin(\theta) = \dfrac{\text{opposé}}{\text{hypoténuse}}$
\end{CodeTeX}
\end{frame}

\begin{frame}[fragile]{Solution de l'exercice « Réparer »}
  \begin{columns}[T]
    \begin{column}{0.48\textwidth}
\begin{CodeTeX}[title={Code corrigé}]
\section{Exercice}
Calculer $x^2$ si $x = 3$.
\begin{itemize}
  \item La fraction $\frac{1}{2}$.
  \item Le rabais : 20\%.
\end{itemize}
\end{CodeTeX}
    \end{column}
    \begin{column}{0.48\textwidth}
      \small
      \begin{enumerate}
        \item Accolade fermante oubliée après le titre de section.
        \item \texttt{x\^{}2} hors du mode mathématique.
        \item Accolade fermante oubliée dans \cmd{frac}.
        \item \texttt{\%} non protégé : la fin de la ligne devient un
          commentaire.
        \item \cmd{begin\{itemize\}} fermé par
          \cmd{end\{enumerate\}}.
      \end{enumerate}
    \end{column}
  \end{columns}
\end{frame}
